\chapter{Limits as the number of subsequent predictions increases}
\label{sec:asymptotic-limits}

In this chapter the number \(r\geq1\) of fixed initial scores, their common
standardized value \(a\), and the correlation \(c\) stay fixed. The number
\(k\) of subsequent predictions increases, and \(N=r+k\) is the total
number of values in the inclusive mean. The limits concern the proved
prediction, mean and SD formulas, with the same counting convention as in
Theorems~\ref{thm:multiple-mean} and \ref{thm:multiple-sd}.

\begin{proposition}[Correlation domain for an unbounded sequence of counts]
\label{prop:limit-domain}
For fixed \(r\), the covariance structure is admissible for arbitrarily
large finite \(k\) exactly when \(0\leq c\leq1\).
\end{proposition}
\begin{proof}
Proposition~\ref{prop:domain}, with total size \(N=r+k\), requires
\(1+(N-1)c\geq0\) and \(c\leq1\). A fixed negative \(c\) makes the
first quantity negative for sufficiently large \(N\), so such a
correlation cannot be retained as \(k\) increases without bound.
Conversely, \(0\leq c\leq1\) satisfies both eigenvalue conditions for
every finite \(N\); the Gaussian covariance construction in that
proposition supplies each finite model. This argument concerns
admissibility at all finite sizes and does not require an additional
infinite-process assertion.
\end{proof}

\begin{lemma}[A square-root limit bound]
\label{lem:square-root-limit-bound}
For nonnegative real \(x,y\),
\[
 \lvert\sqrt{x}-\sqrt{y}\rvert\leq\sqrt{\lvert x-y\rvert}.
\]
Consequently, convergence of nonnegative variances implies convergence of
their nonnegative square roots to the square root of the limiting variance.
\end{lemma}
\begin{proof}
Put \(u=\sqrt{x}\) and \(v=\sqrt{y}\). If \(u\geq v\geq0\), then
\((u-v)^2\leq(u-v)(u+v)=x-y\); otherwise interchange \(x\) and \(y\).
Taking nonnegative square roots proves the inequality. If
\(\lvert x_k-y\rvert<\varepsilon^2\), it gives
\(\lvert\sqrt{x_k}-\sqrt y\rvert<\varepsilon\), proving the stated
limit implication directly from the definition of convergence.
\end{proof}

\begin{theorem}[Limits of predictions, inclusive means and actual-mean SDs]
\label{thm:prediction-limits}
Suppose \(0\leq c\leq1\), \(r\geq1\), and put
\(d_r=1+(r-1)c\). Every subsequent prediction is already the constant
\[
 b_r=\frac{rca}{d_r}.
\]
Its limit and the inclusive recursive-mean limit are therefore
\begin{align}
 \lim_{j\to\infty}e_j^{(r)}&=b_r,
     \label{eq:limit-next-standardized}\\
 \lim_{k\to\infty}B_{r,k}&=b_r.
     \label{eq:limit-inclusive-standardized}
\end{align}
For the actual mean \(C_{r,k}\), the population and conditional SD limits
are
\begin{align}
 \lim_{k\to\infty}\operatorname{SD}(C_{r,k})
   &=\sqrt c,
     \label{eq:limit-population-sd}\\
 \lim_{k\to\infty}\operatorname{SD}(C_{r,k}\given Z=a\one_r)
   &=\sqrt{\frac{c(1-c)}{d_r}}.
     \label{eq:limit-conditional-sd}
\end{align}
The fixed recursively filled mean has SD zero for every finite \(k\) and
in the limit.
\end{theorem}
\begin{proof}
Theorem~\ref{thm:multiple-mean} makes \(e_j^{(r)}=b_r\) at every step;
the value of the next prediction does not evolve towards a different
limit. Its inclusive-mean formula can be rearranged exactly as
\begin{equation}
 B_{r,k}
  =b_r+\frac{r(a-b_r)}{r+k}
  =b_r+\frac{ra(1-c)}{d_r(r+k)}.
  \label{eq:mean-limit-error}
\end{equation}
Here \(d_r\geq1\). The error is at most
\(r|a|(1-c)/(d_r k)\); given \(\varepsilon>0\), choosing
\(k>r|a|(1-c)/(d_r\varepsilon)\) makes it smaller than
\(\varepsilon\). If that numerator is zero, the error is identically
zero. This proves the mean limit.

Theorem~\ref{thm:multiple-sd} gives population variance
\(c+(1-c)/N\). Its difference from \(c\) is
\((1-c)/N\), which tends to zero. For the conditional variance, substitute
\(k=N-r\) in \eqref{eq:multiple-conditional-variance} to obtain
\begin{equation}
 \Var(C_{r,k}\given Z=a\one_r)
 =\frac{1-c}{d_r}
  \left[c+\frac{1-(r+1)c}{N}-\frac{r(1-c)}{N^2}\right].
 \label{eq:conditional-variance-limit-expansion}
\end{equation}
Its absolute difference from \(c(1-c)/d_r\) is bounded by
\[
 \frac{1-c}{d_r N}
   \bigl(|1-(r+1)c|+r(1-c)\bigr),
\]
since \(N\geq1\). The right-hand side tends to zero, by the same
explicit \(1/N\) argument as for the mean. All these variances are
nonnegative, so Lemma~\ref{lem:square-root-limit-bound} proves the two SD
limits. The zero SD of the fixed recursive mean follows from
Theorem~\ref{thm:moments} at every count. At \(c=1\), the equal-score
support argument in Proposition~\ref{prop:multiple-boundaries} supplies
the same constant predictions and conditional SD zero directly.
\end{proof}

\begin{proposition}[A coupling of the Gaussian aggregate laws]
\label{prop:gaussian-limit-coupling}
The population or conditional aggregate laws with means \(m_k\) and
variances \(v_k\) considered above admit coupled realizations converging
in mean square to a Gaussian law with the calculated limiting mean
\(m_\infty\) and variance \(v_\infty\), including zero variances.
\end{proposition}
\begin{proof}
Use a single standard normal \(G\), and put
\(Y_k=m_k+\sqrt{v_k}G\) and
\(Y_\infty=m_\infty+\sqrt{v_\infty}G\).
Theorem~\ref{thm:gaussian} gives the required Gaussian marginal laws.
The moment identities and \(\E[G]=0\), \(\E[G^2]=1\) give
\[
 \E[(Y_k-Y_\infty)^2]
 =(m_k-m_\infty)^2+
   (\sqrt{v_k}-\sqrt{v_\infty})^2\longrightarrow0,
\]
because Theorem~\ref{thm:prediction-limits} proves both parameter
limits, using Lemma~\ref{lem:square-root-limit-bound} for the SDs.
This is convergence in mean square by definition. This coupling concerns
representations of the aggregate laws; it does not assert that the
original expanding test battery uses this single-normal construction.
\end{proof}

\begin{corollary}[Limits for the IQ-150 examples]
\label{cor:iq150-limits}
On the reporting scale \(Q=100+15X\), with initial IQ scores \(150\),
\(a=10/3\), and selected correlation \(c=3/5\), the limiting inclusive
mean is also the IQ prediction for every subsequent test:
\begin{equation}
 m_{r,\infty}^{\mathrm{IQ}}=q_r
 =100+\frac{150r}{3r+2}.
 \label{eq:iq150-limit-means}
\end{equation}
The population SD limit is \(\sqrt{135}\) IQ points for all three initial
counts. The conditional SD limits for \(r=1,2,3\) are respectively
\begin{equation}
 \sqrt{54},\qquad \sqrt{135/4},\qquad \sqrt{270/11}
 \quad\text{IQ points}.
 \label{eq:iq150-limit-sds}
\end{equation}
\end{corollary}
\begin{proof}
Corollary~\ref{cor:iq-reporting} and Theorem~\ref{thm:moments} transform a
standardized mean by \(100+15b_r\), and an SD by multiplication by
\(15\). Substituting \(c=3/5\) gives
\(d_r=(3r+2)/5\) and \(100+15b_r=100+150r/(3r+2)\).
The limiting population variance on this scale is \(225c=135\).
The limiting conditional variance is
\(225c(1-c)/d_r=270/(3r+2)\), which gives the three displayed fractions.
\end{proof}

The following table adds the \(k\to\infty\) case to the finite IQ-150
examples. All means include the initial scores. Values are rounded to six
decimal places.
These infinite-count rows are calculated analytically. Simulations at
finite counts can be compared with the finite formulas and the convergence
bounds proved above.
\begin{center}
\begin{tabular}{r r r r}
\hline
Initial scores \(r\) & Limiting mean IQ & Population SD & Conditional SD \\
\hline
1 & \(130.000000\) & \(11.618950\) & \(7.348469\) \\
2 & \(137.500000\) & \(11.618950\) & \(5.809475\) \\
3 & \(140.909091\) & \(11.618950\) & \(4.954337\) \\
\hline
\end{tabular}
\end{center}

The means \(130\) and \(275/2\) are exact terminating values;
\(1550/11\) and its six-decimal representation already have a certificate
in \Cref{sec:certificates}. For completeness, the new SD roundings
have the following exact positive-integer certificates in the notation of
Lemma~\ref{lem:rounding}, with \(T=10^6\):
\begin{center}
\begin{tabular}{l r r r}
\hline
Quantity & Exact variance & \(L_T\) & \(U_T\) \\
\hline
Population limit & \(135\) & 50065799 & 42885801 \\
Conditional limit, \(r=1\) & \(54\) & 42818031 & 15969721 \\
Conditional limit, \(r=2\) & \(135/4\) & 96541596 & 89361604 \\
Conditional limit, \(r=3\) & \(270/11\) & 193169781 & 242811875 \\
\hline
\end{tabular}
\end{center}
Each integer is positive, so the lemma proves each displayed SD rounding.


\section{Two distinct large-count limits: initially observed versus future tests}
\label{sec:two-count-limits}
The preceding limits hold the number of initial observations \(r\) fixed
and let the subsequent count \(k\) grow. The arbitrary-initial-count
result in Corollary~\ref{cor:many-equal-starts} instead holds \(k\)
fixed and lets the number \(n=r\) of initially observed scores equal to
\(a\) grow. Interchanging these operations requires separate analysis.

\begin{theorem}[Order of the two sequential-count limits]
\label{thm:two-count-limits}
For fixed \(a\in\mathbb R\) and \(0\leq c\leq1\), let
\[
 B_{n,k}=\frac{n[1+(n+k-1)c]a}
                  {(n+k)[1+(n-1)c]},\qquad n,k\geq1.
\]
When \(c=0\), the iterated limits generally differ:
\[
 \lim_{n\to\infty}\lim_{k\to\infty}B_{n,k}=0,\qquad
 \lim_{k\to\infty}\lim_{n\to\infty}B_{n,k}=a.
\]
They coincide only if \(a=0\). For every fixed \(0<c\leq1\),
both iterated limits equal \(a\). In contrast, the population SD has
both iterated limits \(\sqrt c\), and the conditional SD of the
actual inclusive mean has both iterated limits zero, throughout
\(0\leq c\leq1\).
\end{theorem}
\begin{proof}
At \(c=0\), the general mean simplifies to \(na/(n+k)\).
For each fixed \(n\), it tends to zero as \(k\) increases; for each
fixed \(k\), it tends to \(a\) as \(n\) increases. These yield the
two stated orders of limits. For fixed \(c>0\), the limit as \(k\)
grows is \(b_n=nca/[1+(n-1)c]\), by
Theorem~\ref{thm:prediction-limits}; dividing numerator and denominator
by \(n\) shows \(\lim_{n\to\infty}b_n=a\). In the reverse order,
Corollary~\ref{cor:many-equal-starts} proves
\(\lim_{n\to\infty}B_{n,k}=a\) for each fixed \(k\), and the
outer \(k\)-limit leaves this constant unchanged.

The population variance
\(\operatorname{Var}(C_{n,k})=c+(1-c)/(n+k)\)
has limit \(c\) in either count order. For conditional variance, the
fixed-\(k\), growing-\(n\) limit is zero by
Corollary~\ref{cor:many-equal-starts}, whereas the fixed-\(n\),
growing-\(k\) limit is \(c(1-c)/[1+(n-1)c]\)
by Theorem~\ref{thm:prediction-limits}. The latter tends to
zero as \(n\) grows if \(c>0\), and is identically zero if
\(c=0\) or \(c=1\). Continuity of the nonnegative square root,
proved in Lemma~\ref{lem:square-root-limit-bound}, transfers these
variance limits to the corresponding SD limits. Fixed negative
correlations are excluded, since the covariance matrix cannot remain
positive semidefinite along either unbounded sequence of counts.
\end{proof}
